\chapter{Bindings and Behavior Scopes}
\label{ch:bindings-behaviors}
\chapterquote{A binding is a small promise: this component and this model will mirror each other until the owner closes the promise.}{The binding rule}

\begin{chaptermap}
Core question & How should Swing components be connected to Corusco models without unowned listener code?\\
Primary APIs & \api{Binding}, \api{BindingScope}, \api{BindingFactory}, \api{BehaviorScope}, \api{ViewBehavior}, \api{BehaviorDescriptor}, \api{BehaviorPhase}, \api{BehaviorCardinality}, \api{StandardBehaviors}, and \api{SwingEdt}.\\
Plain Swing baseline & Component listeners, document listeners, action maps, tooltips, borders, and focus behavior are ordinary Swing machinery; bindings give the machinery an owner.\\
Corusco connection & Direct bindings handle simple component/model links; behaviors handle repeatable component capabilities and generated plans.\\
Companion example & \swingIntegrationExample.\\
\end{chaptermap}

\section{A binding is owned synchronization}

\termdef{Binding}{in Corusco Swing} is a disposable relationship between Swing
state and model state. It usually installs component listeners, subscribes to
model values, initializes component state, guards against feedback loops, and
removes its work when closed. A binding is not merely a helper method that
copies one value once. It is a live relationship with a lifetime.

Plain Swing has all the necessary primitives. A text field has a document. A
document accepts listeners. A model has values. A label can be updated. The
problem is ownership. Who removes the document listener? Who removes the value
subscription? Who prevents a model update from echoing back as a user edit? A
binding answers those questions.

\begin{lstlisting}[caption={A direct text-field binding.}]
Binding binding = BindingFactory.textField(nameField, form.name);

// Later, when the view closes:
binding.close();
\end{lstlisting}

The binding does not own \api{nameField} and does not own \api{form.name}. It
owns the relationship it installed between them.

\begin{principlebox}{Relationship Ownership}
If code installs synchronization between a component and a model, the same code
path should expose a close handle for that synchronization.
\end{principlebox}

\begin{figure}[htbp]
\centering
\begin{minipage}[t]{0.48\linewidth}
\includegraphics[width=\linewidth]{\bindingFormFigure}
\end{minipage}\hfill
\begin{minipage}[t]{0.48\linewidth}
\includegraphics[width=\linewidth]{\behaviorPlanFigure}
\end{minipage}
\caption{Bindings and behaviors are visible relationships over ordinary Swing components: a form mirrors model state, and a generated plan installs named capabilities into one lifecycle scope.}
\label{fig:bindings-behavior-screens}
\end{figure}

A binding is easiest to misunderstand when it is introduced as less code. It is
less handwritten listener code, but that is not its main purpose. Its main
purpose is to make a relationship explicit enough that a reviewer can ask
ordinary engineering questions about it. Which side initializes first? Which
events are considered user edits? Which events are considered model refreshes?
How is feedback suppressed? What state is restored on close? Which objects are
owned by the binding, and which objects merely participate in it? A helper that
copies text from a field into a model answers almost none of these questions. A
binding should answer all of them.

Initialization is the first contract. When \api{BindingFactory.textField} is
installed, the component is initialized from \api{TextFieldModel.rawText}. This
direction is not arbitrary. The form model is the durable screen state. The
text field is a Swing presentation of that state. If the field happened to
contain old designer text, previous dialog contents, or a value left by a test,
the binding should not treat that text as a fresh user edit merely because the
listener has just been installed. The model speaks first; then user edits are
observed.

This rule is especially important in dialogs and generated forms. A dialog may
reuse a panel class for several records. A generated form model may be built
from a record snapshot. The visual components may be constructed before the
model is available, because layout code wants concrete components. Binding
installation is the moment when the model and the components are joined. If the
join direction is vague, initialization bugs look like validation bugs, dirty
tracking bugs, or mysterious "field changed during open" bugs.

The second contract is origin. Corusco value changes carry enough information
to distinguish user-originated updates from programmatic updates when the model
supports it. Text and selected-state bindings write with
\api{StandardChangeOrigin.USER}. A presenter can then distinguish a user's keystroke
from a reset button, a background load, or an apply operation. That distinction
is not decorative. Dirty flags, validation timing, undo grouping, command
enablement, and support diagnostics all depend on knowing whether a change came
from the person at the keyboard or from application code.

Plain Swing listeners frequently lose this fact. A document listener observes a
document change, but the document does not know whether the application called
\api{setText} or the user typed a character. A careless listener therefore
marks a form dirty while a model is being loaded, or fires validation while a
field is being initialized. Binding code cannot make every semantic decision
for the application, but it can preserve the origin at the boundary where Swing
events become Corusco model changes.

The third contract is feedback suppression. A two-way relationship has a simple
failure mode: model update sets component text; component document event writes
the same text back to the model; model emits another change; the component is
set again. Some loops terminate because the value is equal. Some do not,
especially when formatting, parsing, trimming, or validation changes the text.
Even a terminating loop can produce duplicate events and confusing origin
metadata. Corusco bindings use a feedback guard around programmatic component
updates so component listeners do not reinterpret them as user edits.

This is one reason the binding should be narrow. If a text binding also tries
to own a border, tooltip, status label, dirty marker, and command state, the
feedback story becomes harder to inspect. A direct text binding should
synchronize text. A validation tooltip binding should own tooltip text. A
validation border binding should own border decoration. Narrow relationships
make the feedback guard local and make cleanup obvious. Rich screens are built
by composing narrow bindings, not by writing one omniscient listener.

The fourth contract is cleanup. Closing a binding removes the listeners and
subscriptions that the binding installed. It does not dispose the text field.
It does not dispose the form model. It does not clear the resource bundle or
release the presenter. The distinction matters because Swing views have
lifetimes that are shorter than application models. A dialog may close while
the model snapshot remains useful for tests. A tool window may be hidden and
shown again. A generated plan may be installed for one view activation and
closed before the next. Cleanup must detach the relationship without pretending
to own the participants.

The fifth contract is restoration. Some bindings are pure subscriptions and do
not restore final component state. \api{BindingFactory.labelText} removes the
value subscription when closed, but it does not restore the label's old text.
That is usually the right contract: the label was a view of current state, and
its final text is harmless after the view is closing. Other bindings
temporarily claim a component property and must restore it. A validation border
binding snapshots the original border and restores it on close. A tooltip
binding snapshots the original tooltip slot and restores it on close. A status
binding restores the shared status label when the focus relationship ends.

These policies should be written down because restoration bugs often survive
manual testing. A dialog opens, shows an error border, closes, opens again, and
the new instance inherits a border that belonged to the previous activation. Or
a help tooltip is installed, the presenter is replaced, and the old help text
remains. A binding that changes component state temporarily should have a
visible restoration rule. A binding that intentionally leaves the component in
the current state should say so as well.

The sixth contract is the dispatcher boundary. A core \api{ReadableValue} does
not become Swing-aware because a Swing component subscribes to it. The binding
factory methods require the EDT because they install listeners and mutate
components. If a background task produces model changes that are observed by
Swing bindings, an EDT adaptation must sit between worker completion and Swing
state. The failure should happen near installation or update delivery, not as a
random component exception later.

This is also a documentation virtue. When a binding requires the EDT, examples
should show it being installed from view construction, presenter activation, a
dialog factory, or an EDT-aware test harness. They should not imply that any
thread may bind a component because the underlying core value is thread-neutral.
The core package and the Swing package cooperate precisely because each keeps
its own boundary honest.

Finally, a binding should be boring to close. Calling \api{close} twice should
not remove different things the second time. Adding a binding to a closed scope
should not leave a relationship half-installed. Closing in reverse order should
handle dependencies where later visual decorations depend on earlier primary
bindings. These are lifecycle details that handwritten listener code rarely
addresses until a screen is reopened many times in a test. Corusco's binding
scope exists so those details become ordinary, not heroic.

\textbf{EDT confinement.}

Corusco Swing binding factories require the Event Dispatch Thread. They install
Swing listeners and mutate Swing components. Core values do not dispatch to the
EDT automatically. Therefore the boundary is explicit: call binding installation
on the EDT, update Swing components on the EDT, and use an EDT adapter or task
callback when background work feeds a bound model.

\api{SwingEdt.requireEdt} is deliberately strict. Failing fast is better than
quietly installing a binding that sometimes updates a component from a worker
thread.

\begin{pitfallbox}{The Core/Swing Split}
A \api{ReadableValue} is toolkit-neutral. A \api{BindingFactory} method is
Swing-specific. Do not expect the former to enforce the rules of the latter.
\end{pitfallbox}

This split is the organizing idea of Corusco Swing. Core models describe screen
facts without opening a window: raw field text, parsed values, dirty state,
problem sets, command state, row lists, and table descriptors. Swing bindings
then adapt those facts to concrete components. A \api{JTextField} can display
raw text, but the form model owns the raw text. A \api{JButton} can display an
operation, but the command owns the operation identity. A \api{JTable} can show
rows, but the table descriptor owns column meaning. Once this separation is
understood, bindings stop looking like a convenience wrapper and start looking
like the boundary where the two worlds meet.

The boundary is deliberately explicit. If a background task changes a core
value and that value is observed by a Swing binding, the task still needs a
Swing-aware callback path or an EDT adapter. This may feel stricter than some
listener code written in old applications, but the strictness pays for itself.
The bug is found where the boundary is crossed, not later as a random repaint,
document, or selection failure.

\textbf{Direct binding contracts.}

\api{BindingFactory.textField} binds a \api{JTextField} to a
\api{TextFieldModel}. It initializes the component from raw text, installs a
document listener, propagates user document changes to the model with
\api{StandardChangeOrigin.USER}, observes model raw-text changes, and closes both
listener and subscription when closed.

The phrase "raw text" matters. The binding does not bind the field directly to
the semantic value. It binds Swing text to \api{TextFieldModel.rawText}. Parsing
and semantic value protection remain in the form model.

\begin{lstlisting}[caption={Text binding preserves invalid intermediate input.}]
TextFieldModel<CustomerEdit, BigDecimal> limit = form.creditLimit;
Binding binding = BindingFactory.textField(limitField, limit);

limitField.setText("12,");       // User-originated document change in real UI.
ProblemSet parseProblems = limit.problems();
\end{lstlisting}

The listener code is not gone. It is inside a tested binding with feedback
guards and cleanup.

\textbf{Selected-state binding.}

\api{BindingFactory.selected} binds an \api{AbstractButton} to a Boolean
\api{FieldModel}. This covers check boxes, radio buttons, and toggle buttons.
The component is initialized from the model. User item events write
\api{StandardChangeOrigin.USER}. Model events update the button while a guard prevents
echo writes.

\begin{lstlisting}[caption={Checkbox state belongs to the model.}]
FieldModel<CustomerEdit, Boolean> active = form.active;
Binding binding = BindingFactory.selected(activeCheckBox, active);
\end{lstlisting}

This is especially valuable when the same Boolean is also represented by a menu
item or command. One model owns the fact; controls present it.

\textbf{Label and status bindings.}

Some bindings are one-way. \api{BindingFactory.labelText} initializes a label
from a readable string value and updates it when the value changes. It does not
restore the previous label text on close. \api{BindingFactory.statusText}
temporarily publishes status text while a component owns focus and restores the
previous status on focus loss or close.

These differences are part of the binding contract. A binding should document
what it restores and what it leaves as final state. Without that documentation,
closing a view becomes guesswork.

\textbf{Validation decorations.}

Validation feedback is not the same as value synchronization. A text binding
mirrors text. A validation tooltip binding observes problems and writes a
tooltip. A validation border binding observes problems and decorates a component
border. These concerns should be separate so they can be combined, tested, and
closed independently.

\begin{lstlisting}[caption={Separate binding concerns for one field.}]
scope.add(BindingFactory.textField(nameField, form.name));
scope.add(BindingFactory.validationTooltip(nameField, form.name.problemSet()));
scope.add(BindingFactory.validationBorder(nameField, form.name.problemSet()));
\end{lstlisting}

The component has one visual surface, but several relationships contribute to
it. When this pattern repeats across generated fields, direct manual binding can
become verbose. That is where behaviors enter.

\textbf{Binding scopes.}

\api{BindingScope} owns bindings for one view or dialog lifecycle. It is the
Swing-specific counterpart to the core subscription scope. Add each binding as
the view is assembled; close the scope when the view closes.

\begin{lstlisting}[caption={A direct binding scope.}]
final class CustomerFormPanel extends JPanel implements AutoCloseable {
    private final BindingScope bindings = new BindingScope();

    CustomerFormPanel(CustomerEditForm form) {
        bindings.add(BindingFactory.textField(nameField, form.name));
        bindings.add(BindingFactory.selected(activeBox, form.active));
        bindings.add(BindingFactory.labelText(statusLabel, form.statusText()));
    }

    @Override
    public void close() {
        bindings.close();
    }
}
\end{lstlisting}

This is already a significant improvement over listener sprawl. But it still
requires handwritten installation order and duplicate patterns for every field.

\section{Direct bindings in real screens}

Direct bindings are the right starting point because they keep the machinery
visible. A moderate beginner can open a panel constructor, see a
\api{BindingScope}, see a text field binding, see a selected-state binding,
and see the scope close with the panel. Nothing has been generated yet. No
behavior descriptor has to be understood before the first benefit appears. The
screen is still ordinary Swing: \api{JTextField}, \api{JCheckBox},
\api{JLabel}, \api{JPanel}, MigLayout, and listeners under the surface. The
difference is that the listener relationships now have names and lifetimes.

The first real-screen rule is to install bindings after static component setup
and before user interaction. Set columns, renderers, input verifiers, component
names, accessible static text, and layout constraints before binding when those
properties are construction facts. Then install the binding and let it own the
properties it promises to own. If code calls \api{nameField.setText} after
\api{BindingFactory.textField} has been installed, the write may be a genuine
programmatic model operation expressed through Swing, or it may be an accidental
fight with the model. Prefer changing the model when the bound fact changes.
Use direct component writes after binding only when the binding contract says
they are part of the relationship.

Text fields are the clearest example. A text component has a \api{Document};
the document has listeners; the form model has raw text, parse state, semantic
value, and problems. The binding joins the document to raw text. It does not
make parsing a Swing concern. It does not require every intermediate text to be
a valid semantic value. This is why a user can type the beginning of a decimal
or date without the component immediately erasing the text. Raw text belongs to
the editing experience; semantic value belongs to the form model once parsing
and validation allow it.

This distinction is an important bridge from Swing to Corusco. A beginner may
ask why a binding is not simply \api{JTextField} to \api{BigDecimal}. The answer
is that editing is not assignment. The string \api{"12,"} may be a perfectly
reasonable intermediate text in a locale-aware decimal field and still not be a
valid \api{BigDecimal}. A field model can remember the raw text, keep the last
valid semantic value, and expose a problem. A direct Swing listener that writes
straight to a \api{BigDecimal} property either rejects intermediate text too
early or invents error state beside the model. Binding to raw text keeps the
editing boundary explicit.

Text areas follow the same rule but increase the importance of ownership.
Large text areas may participate in document undo, line wrapping, scroll panes,
and validation summaries. A text-area binding should still be narrow: document
to raw text. Word count, dirty marker, save command, and validation decoration
can observe model facts or document-derived values separately. This keeps the
text area from becoming a private model hidden in a widget.

Selected-state bindings look simpler, but they teach the same lesson. A
checkbox's selected state mirrors a Boolean field. User item events write with
\api{StandardChangeOrigin.USER}; model changes update the checkbox under a feedback
guard. If the same Boolean also appears as a menu item, toolbar toggle, or
filter chip, the Boolean field remains the fact. The controls are presentations
of it. Plain Swing often lets each control become a little source of truth,
with synchronization listeners trying to keep the controls equal. Direct
bindings reverse the direction: one model fact, many possible presentations.

Button enabled-state binding is intentionally one-way. A command or readable
Boolean says whether the button is enabled. The user does not edit the enabled
state by clicking a disabled button. This sounds obvious, but it is a good
review test. A two-way binding should exist only when the user can directly edit
the model fact through the component. A one-way binding is the right tool when
the component presents a derived fact. Overusing two-way binding is another
form of hidden ownership confusion.

Labels are also one-way. A status label, total count label, selected-customer
label, or validation summary label presents a readable value. Closing the
binding removes the subscription. It does not need to restore old text unless
the label is shared and temporarily claimed. This difference between an owned
label and a shared status label is small but practical. A panel-owned label dies
with the panel. A status label in a shell may outlive the field that borrowed
it while focused. The binding's restoration policy follows the lifetime of the
component property it claims.

Status bindings show how focus state enters the relationship. The component
does not permanently own the status bar. It borrows it while focused, snapshots
the old text, publishes its status, follows status value changes while active,
and restores the snapshot when focus leaves or the binding closes. That is much
harder to implement correctly with one focus listener per field and no close
handle. It is also the kind of detail that users notice: status text that stays
behind after focus moves makes the screen feel unreliable.

Validation decoration should be composed from facts, not duplicated from parse
logic. A field model exposes problems. A tooltip binding converts those
problems into tooltip text according to policy. A border binding converts
problem severity into component border state. A summary panel may observe the
same problem set and render a list. The binding chapter must insist on this
separation because otherwise validation becomes invisible listener code again.
If a tooltip listener recomputes validation by reading the text field, it is
not a tooltip binding; it is a second validation engine hiding in Swing code.

Tooltip ownership deserves special care because Swing gives each component one
tooltip text slot. Validation, disabled reason, static field help, and F1 help
indicator all want to explain the same component. If each concern writes the
slot independently, the last writer wins and the losing explanations disappear.
\api{BindingFactory.composedTooltip} and the corresponding standard behavior
exist so the slot has one owner and several contributors. The result is not
only better text; it is better lifecycle. The original tooltip can be restored
on close because one binding owns the slot.

Borders have a similar but more visual conflict. A Look-and-Feel may install a
border that encodes focus, error state, disabled state, or theme density. A
validation border binding that simply replaces it with a red line may be
acceptable as a minimal implementation, but production decoration often needs a
composed border or a layer-based mark that respects the original border. The
important binding contract remains: snapshot what you replace, update only what
you own, and restore on close. Later chapters on dialogs and validation can
refine the visual policy; this chapter establishes the ownership rule.

Accessible text is not decoration in the trivial sense. It is secondary to the
primary value binding, but it is part of the component's public contract. A
binding or behavior can set accessible name and description from field
metadata, resources, or explicit text, and restore the previous values on
close. Doing this through the same scope as other component relationships has a
useful side effect: tests can construct a field, install the plan, and assert
that visible text, tooltip, validation, and accessible description agree. The
screen becomes easier to audit for both visual and nonvisual users.

The direct-binding assembly order can be described as a small ritual. Build
models. Build components. Apply static component configuration. Lay out the
components. Create a scope. Add primary value bindings. Add one-way
presentation bindings. Add validation and tooltip bindings. Add focus or status
bindings. Store the scope in the owning view or presenter. Close the scope when
the view closes. The ritual is not bureaucracy; it prevents the two common
failures of Swing wiring: hidden listeners that never die and properties that
have several owners.

\begin{lstlisting}[caption={Direct binding assembly order for a small form.}]
final class CustomerSummaryPanel extends JPanel implements AutoCloseable {
    private final BindingScope bindings = new BindingScope();

    CustomerSummaryPanel(CustomerEditForm form, JLabel shellStatus) {
        JTextField name = new JTextField(24);
        JCheckBox active = new JCheckBox();
        JLabel summary = new JLabel();

        setLayout(new MigLayout("fillx", "[][grow]", "[][][]"));
        add(new JLabel("Name"));
        add(name, "growx, wrap");
        add(new JLabel("Active"));
        add(active, "wrap");
        add(summary, "span 2, growx");

        bindings.add(BindingFactory.textField(name, form.name));
        bindings.add(BindingFactory.selected(active, form.active));
        bindings.add(BindingFactory.labelText(summary, form.summaryText()));
        bindings.add(BindingFactory.composedTooltip(
                name, form.name.problemSet(), null,
                "Customer display name", true));
        bindings.add(BindingFactory.validationBorder(name, form.name.problemSet()));
        bindings.add(BindingFactory.statusText(name, shellStatus, "Edit customer name"));
    }

    @Override
    public void close() {
        bindings.close();
    }
}
\end{lstlisting}

The listing is deliberately direct. It does not use a behavior scope because
the reader should first see the relationships plainly. The same form could be
assembled by generated behavior plans later, but the generated plan should not
hide any concept that is absent here. It should merely remove repetition and
add conflict metadata.

Direct bindings also help migration because they have a small blast radius.
Pick one field with a document listener. Replace only the document listener
with \api{BindingFactory.textField}. Keep the rest of the screen unchanged.
Add the binding to a scope and close the scope where the old listener would
have been removed if the old code had supported removal. Then replace tooltip
and border wiring. Then replace status text. This path lets a large legacy
screen become more disciplined without a rewrite.

During migration, look for writes that cross the binding boundary. If old code
still calls \api{nameField.setText} from a presenter after the field is bound,
ask whether it should instead call \api{form.name.setRawText} with model
origin. If old code reads \api{nameField.getText} during save, ask whether it
should read the form model's parsed value or raw text. If old code inspects a
tooltip to decide whether a field is valid, ask why it is not reading the
problem set. Each of these questions moves state out of widgets and back into
models.

Direct bindings should be tested by behavior, not by implementation trivia.
Install the binding on the EDT. Assert initial component state. Simulate a
user-like component change and assert the model. Change the model and assert
the component. Close the binding or scope. Change one side again and assert the
other side no longer follows. For a tooltip or border binding, assert both the
active state and the restored state. These tests are small, repetitive, and
valuable because they guard the lifecycle contract rather than a single visual
snapshot.

The direct-binding style has a natural limit. When five fields all need the
same text binding, validation tooltip, validation border, accessible text, and
select-all-on-focus behavior, handwritten direct bindings become a recipe
copied many times. Copying is tolerable for two fields and dangerous for
twenty. One field forgets restoration. Another gets tooltip before validation.
A third misses accessible text. The answer is not to invent a larger helper
method that hides everything. The answer is a behavior plan: keep each
capability named, but let metadata and generated code assemble the repeated
list.

\section{Behavior scopes and generated plans}

\termdef{View behavior}{in Corusco Swing} is a named, ordered, lifecycle-owned
component capability. Text binding, selected-state binding, validation tooltip,
validation border, composed tooltip, select-all-on-focus, accessible text,
help-on-F1, command binding, and busy overlay can all be behaviors.

A behavior does not touch a component when it is created. It becomes active when
a \api{BehaviorScope} installs it with a \api{BehaviorContext}. Installation
returns a \api{Binding}; the scope owns that binding.

\begin{lstlisting}[caption={Installing behaviors rather than individual bindings.}]
BehaviorScope scope = new BehaviorScope(helpService);

scope.install(nameField, List.of(
        StandardBehaviors.textFieldBinding(form.name),
        StandardBehaviors.composedTooltip(
                form.name.problemSet(), null, "Customer display name", true),
        StandardBehaviors.validationBorder(form.name.problemSet()),
        StandardBehaviors.selectAllOnFocus()));
\end{lstlisting}

The list is longer than a single binding call because the example is explicit.
Generated plans can produce the list from annotations and descriptors.

A behavior is not a magic listener. It is a component responsibility with
metadata. That metadata is the difference between a plan and a pile of helper
objects. Before installation, the behavior can say: this is my key; this is my
phase; I may appear once or many times; I do or do not occupy the component's
primary binding slot. The scope can then reason about the list before the list
quietly changes a live Swing component. This is a modest idea, but it is the
reason generated view plans can be safe enough to use in ordinary code.

The key should describe a responsibility rather than an object instance. The
built-in \api{StandardBehaviorKeys.TEXT\_BINDING} key means "primary text
component binding." It does not mean "the particular object returned by this
factory call." The tooltip key means "the component tooltip slot is owned by
this behavior." The select-all key means "focus gained selects all text." A
stable responsibility key lets tests and diagnostics ask which capabilities are
installed without depending on anonymous class names or generated object
identity.

This has an important review consequence. If a custom behavior's key is
\api{test/my-behavior} or \api{misc/field-hook}, the key is not helping. It
should be named after the component responsibility: \api{interaction/nudge-on-arrow},
\api{decoration/dirty-marker}, \api{binding/date-text}, or
\api{interaction/open-chooser}. Generated plans and tests become readable only
when keys are readable. A behavior key is not user-facing text, but it is
developer-facing documentation.

Phases are deliberately few. \api{BehaviorPhase.BINDING} installs primary value
relationships. \api{BehaviorPhase.DECORATION} installs visual, tooltip, and
accessibility state that decorates the component. \api{BehaviorPhase.INTERACTION}
installs focus, key, help, status, and other user-interaction affordances. More
phases would make ordering look more precise than it really is. The goal is not
to encode every presenter workflow in phase order. The goal is to prevent
primary synchronization, decoration, and interaction from depending on whatever
order an annotation processor, list builder, or refactoring happened to
produce.

The phase model also supports reverse cleanup. If a binding is installed first,
a decoration second, and an interaction third, closing in reverse order removes
the interaction before the decoration and the decoration before the primary
binding. That is usually the least surprising dependency order. A focus status
behavior may depend on the component still existing; a validation decoration
may depend on model state still being available; the primary binding is the
foundation. Reverse close is not a universal proof of correctness, but it is a
sound lifecycle default and matches the subscription-scope discipline used in
the core package.

Cardinality answers a different question from phase. Phase says when a behavior
is installed. Cardinality says whether another behavior with the same key may
coexist on the same component. A tooltip owner should be single because Swing
has one tooltip string slot. A validation border should be single because two
border owners cannot both restore the same original border correctly. A command
or key behavior may be multiple if each instance owns a distinct key, command,
or action slot. The descriptor should express the truth of the component
resource being used.

Primary binding conflict is stricter than ordinary single-cardinality. Two
different behavior keys may still both claim to be primary bindings. A text
binding and a hypothetical formatted-value binding could have different keys
and still be incompatible on the same field because both want to own the main
component/model synchronization. This is why \api{BehaviorDescriptor.primaryBinding}
sets both single cardinality and the primary-binding conflict flag. The scope
tracks primary binding by component, not merely by key.

The distinction prevents a subtle bug in extensible systems. Suppose a generated
plan installs a standard text binding. A handwritten extension later installs a
custom masked-text binding with its own key. Duplicate-key detection would not
catch the conflict because the keys are different. Primary-binding conflict
detection does catch it because both behaviors claim the primary slot. The
failure appears when the screen is assembled, before the user discovers that
typing in the field updates two models or alternates between two formatters.

\begin{lstlisting}[caption={Descriptor choices encode component responsibility.}]
BehaviorDescriptor text =
        BehaviorDescriptor.primaryBinding(StandardBehaviorKeys.TEXT_BINDING);

BehaviorDescriptor tooltip =
        BehaviorDescriptor.single(StandardBehaviorKeys.TOOLTIP,
                BehaviorPhase.DECORATION);

BehaviorDescriptor command =
        BehaviorDescriptor.multiple(BehaviorKey.of("interaction/command"),
                BehaviorPhase.INTERACTION);
\end{lstlisting}

The example is schematic, but the categories matter. The text binding owns the
main data flow. The tooltip owns one Swing slot. The command interaction may be
multiple when several command bindings install different actions. A custom
behavior should be able to justify its descriptor in these terms. If it cannot,
the behavior probably mixes responsibilities.

Behavior context is intentionally narrow. It gives the behavior the target
component, the owning scope, and optional services such as help. It does not
give the behavior arbitrary access to the application. This keeps behavior
installation from becoming a second dependency-injection system. A behavior may
need a help service to open F1 help. It may need a resource-resolved string
captured at factory time. It may need a command object captured by a command
behavior. It should not discover repositories, windows, database sessions, or
presenter workflows by reaching through global state.

The help behavior illustrates the context boundary. The behavior factory
captures a \api{HelpTopic}. During installation, the context supplies the
\api{HelpService} from the owning behavior scope. The behavior then installs an
input-map and action-map entry for F1 and restores the previous entries during
cleanup. The help service is not global, and the target component is not hidden.
Tests can create a scope with a fake service, install the behavior, invoke the
action, and assert the request.

Generated behavior plans should remain direct. The processor-generated
\api{CustomerEditBehaviorPlan} does not invent a runtime reflection engine. It
calls view accessors and installs ordinary standard behaviors. For a text field
it installs text binding, validation tooltip, validation border, and
select-all-on-focus. For a checkbox it installs selected-state binding. For a
radio group it installs option-valued selected-state binding against generated
option descriptors. The generated plan is valuable precisely because it is
boring. A reader can open it and see the same calls a careful developer would
have written by hand.

\begin{lstlisting}[caption={The shape of a generated text-field plan.}]
scope.install(view.nameField(), List.of(
        StandardBehaviors.textFieldBinding(model.name),
        StandardBehaviors.validationTooltip(model.name.problemSet()),
        StandardBehaviors.validationBorder(model.name.problemSet()),
        StandardBehaviors.selectAllOnFocus()
));
\end{lstlisting}

The generated code still leaves layout to the view. This is an important design
choice for a Swing book that recommends MigLayout. A generated form plan should
not decide where fields are placed, how dense the dialog is, which controls
belong in a toolbar, or how a business screen adapts to resizing. The plan owns
repeatable behavior, not geometry. That separation keeps generated code useful
in serious Swing screens instead of turning it into a form designer with the
usual limitations of form designers.

Generated and handwritten behavior can therefore cooperate. The generated plan
installs the standard field capabilities. The handwritten presenter may add
status text, F1 help, a field-specific chooser key, a diagnostic dirty marker,
or a command binding. Because both generated and handwritten behaviors pass
through the same scope, they share conflict detection and cleanup. The custom
addition does not need to know whether the standard behaviors were generated or
written by hand; it only needs to obey the same descriptor contract.

The reverse is also true: generated code cannot protect the screen from
relationships installed outside the scope. If a legacy listener is added
directly to the field and writes another model, the behavior scope cannot see
it. During migration, direct listener code should be treated as outside the
plan and reviewed accordingly. The goal is not to forbid all direct listeners.
Some local Swing interactions are best written directly. The goal is to keep
model synchronization, repeated decoration, and generated capabilities inside a
scope where they can be ordered, inspected, and closed.

Failure during behavior installation should be considered a plan defect. A
duplicate single-cardinality key means two behaviors want the same component
responsibility. A second primary binding means two behaviors want the same
source of truth. A missing help service means a behavior requiring help was
installed into a scope that cannot provide it. A behavior that throws while
installing has not completed its contract. The tests in the codebase deliberately
assert these failures because silent partial installation is worse than early
rejection.

There is one nuance: if a behavior fails during installation after earlier
behaviors were already installed, the scope still owns the earlier returned
bindings. Closing the scope remains necessary. This is normal lifecycle
engineering. A caller that catches an installation failure should close the
scope or let the owning construction failure path close it. Generated plans
should usually fail fast and abort view activation. Handwritten assembly code
should be equally strict unless it is intentionally installing optional
capabilities with local fallback.

The behavior-scope API exposes installed keys for tests. This is not meant to
replace effect assertions. A good test may assert that a generated plan
installed \api{StandardBehaviorKeys.SELECT\_ALL\_ON\_FOCUS}; another test may
simulate focus and assert selected text. The key assertion proves the plan
contains a capability. The effect assertion proves the capability works. Both
are useful at different levels. Generated-code tests often care about the key
set; component behavior tests care about the visible and interactive effect.

The small size of \api{BehaviorPhase}, \api{BehaviorKey}, and
\api{BehaviorDescriptor} is deliberate. They are not a grand UI framework. They
are enough metadata to make assembly deterministic, inspectable, and safe
against common conflicts. Swing itself still supplies components, listeners,
input maps, action maps, tooltips, borders, focus events, and painting. Corusco
supplies ownership and plan structure around those mechanisms.

\textbf{Behavior descriptors.}

\api{BehaviorDescriptor} names a behavior with a \api{BehaviorKey}, phase,
cardinality, and primary-binding conflict policy. This metadata lets a
\api{BehaviorScope} validate a plan before silent double installation.

Without descriptors, installing two text bindings on the same field might
appear to work until events bounce between two models. With descriptors, primary
binding conflict is a defined error.

\begin{coruscobox}{Conflict Detection}
Corusco behavior scopes detect duplicate single-cardinality behaviors and
multiple primary bindings on the same component. The failure happens during
installation, not after a user reports strange text-field behavior.
\end{coruscobox}

\textbf{Phases.}

Behaviors are sorted by \api{BehaviorPhase} before installation. Phases let a
plan say that primary bindings, decorations, help, and other capabilities have
a deterministic order. This is useful for generated code and tests. It also
prevents accidental dependence on list order in annotation processing.

Do not put business policy into phase ordering. Phases are for installation
structure, not for modeling domain workflows.

\textbf{Cardinality.}

Some behaviors may appear more than once. Others should appear at most once per
component. Tooltip behavior is a good example. Swing components have one
tooltip slot, so two independent tooltip writers would conflict. A composed
tooltip behavior should own that slot and combine contributors.

\api{BehaviorCardinality.SINGLE} lets the scope reject duplicate behavior keys.
This is better than letting the last installed behavior win silently.

\textbf{Primary bindings.}

A primary binding is the behavior that owns the main value synchronization for a
component: text field to text model, checkbox to Boolean model, combo box to
selection model, and so on. A component should not normally have two primary
bindings. The behavior scope tracks this and throws a
\api{BehaviorConflictException} when a second primary binding is installed.

\begin{pitfallbox}{Two Masters}
If one text field is bound to two text models, neither model is clearly the
owner. Reject the plan; do not attempt to arbitrate with listener ordering.
\end{pitfallbox}

\textbf{Behavior context.}

\api{BehaviorContext} provides the target component, the owning behavior scope,
and optional services such as \api{HelpService}. Custom behaviors should obtain
application services from context or explicit constructor arguments, not from
globals. This keeps tests and generated plans predictable.

\textbf{Generated behavior plans.}

Generated \api{@CoruscoForm} companions can produce behavior plans for a view
interface. The application still writes the panel and lays it out with
MigLayout. The generated plan installs standard bindings and decorations based
on descriptors.

This division is intentional:

\begin{itemize}
\item The handwritten view owns component construction and layout.
\item The generated plan owns repetitive field behavior.
\item The behavior scope owns installed bindings.
\item The presenter or form model owns state.
\end{itemize}

\begin{migrationbox}{From Generated Panel to Generated Behavior}
Do not require generation to own layout. Let MigLayout and the handwritten view
own geometry. Let generated companions own the repetitive synchronization plan.
\end{migrationbox}

\textbf{Manual and generated behavior can coexist.}

A generated behavior plan should not make handwritten behavior impossible. Real
screens always have a few local decisions: a special status hint, a custom
formatter, a field-specific key binding, or a diagnostic border used only in one
tool window. The behavior scope is designed so generated and handwritten
behaviors can be installed into the same lifecycle, provided they obey the same
descriptor rules.

\begin{lstlisting}[caption={Generated plan plus one local behavior.}]
BehaviorScope scope = new BehaviorScope(helpService);

CustomerEditBehaviorPlan.install(scope, view, form, resources);
scope.install(view.creditLimitField(), List.of(
        new SelectAllOnDoubleClick()));
\end{lstlisting}

The local behavior should not bypass the scope by installing unmanaged
listeners. If it participates in the same component, it should participate in
the same conflict and cleanup policy.

\textbf{Behavior failure modes.}

Behavior installation can fail before a screen is shown. This is preferable to
silent drift. A duplicate single-cardinality tooltip behavior means two writers
want one tooltip slot. A second primary binding means two models want one
component. A missing help service means a help behavior cannot open topics. A
custom behavior that mutates Swing off the EDT violates the scope contract.

Treat these failures as plan errors, not as runtime annoyances. The screen
assembly code should be easy to inspect, and generated plans should have tests
that prove their expected behavior keys.

\textbf{Restoration policy.}

Some bindings restore previous component state when closed; others do not.
Status-text binding restores a shared status label. Validation border binding
should restore the previous border. Label text binding usually does not restore
old text because the final label content may simply be the current state.

Custom behaviors must decide and document restoration policy. Closing a
behavior that temporarily owns a component property should leave the component
in a coherent state. A behavior that permanently initializes a property should
say so.

\begin{center}
\textbf{Binding and behavior review table.}
\end{center}

\noindent\begin{tabularx}{\linewidth}{@{}p{0.30\linewidth}X@{}}
\toprule
Observation & Review question\\
\midrule
Component has two model listeners & Is there one primary binding too many?\\
Tooltip changes unpredictably & Should the contributors be composed by one behavior?\\
Generated plan and local code both bind text & Which binding owns the raw text?\\
Close does not remove a focus listener & Which binding or behavior owns cleanup?\\
Custom behavior reaches a global service & Should the service enter through context?\\
Test inspects private listener arrays & Can installed behavior keys express the assertion?\\
\bottomrule
\end{tabularx}

\textbf{Custom behaviors.}

Write custom behaviors as small reversible units. A custom behavior should
install one capability and return a binding that undoes it. If it installs a
listener, close removes the listener. If it changes a border and promises
restoration, close restores the old border. If it registers a key binding, close
removes or restores the previous mapping.

\begin{lstlisting}[caption={Custom behavior shape.}]
final class SelectAllOnDoubleClick
        implements DecorationBehavior<JTextField> {
    @Override
    public BehaviorDescriptor descriptor() {
        return BehaviorDescriptor.single(
                BehaviorKey.of("select-all-on-double-click"),
                BehaviorPhase.DECORATION);
    }

    @Override
    public Binding install(BehaviorContext<JTextField> context) {
        MouseListener listener = new MouseAdapter() {
            @Override
            public void mouseClicked(MouseEvent event) {
                if (event.getClickCount() == 2) {
                    context.component().selectAll();
                }
            }
        };
        context.component().addMouseListener(listener);
        return () -> context.component().removeMouseListener(listener);
    }
}
\end{lstlisting}

The behavior is small and reversible. It does not know the whole form.

\section{Restoration, composition, and custom reversible behavior}

The hardest part of a reusable Swing behavior is not usually installation. It
is leaving. Installing a focus listener, tooltip, border, key binding, action
map entry, or client property is straightforward. Restoring the previous state
is where careless abstractions leak. A behavior that works in a demo can fail
when a dialog is reopened, when another behavior uses the same component slot,
when Look-and-Feel changes, or when tests install and close the plan several
times in one JVM. A good behavior is therefore designed from its close method
backward.

The close method should undo exactly what installation owned. If installation
adds a listener, close removes that listener and no other listener. If
installation replaces a border, close restores the border that was present at
installation time, unless the behavior participates in a composed border system
with a different documented policy. If installation writes an input-map key,
close restores the previous mapping or removes the mapping when none existed.
If installation subscribes to a model value, close closes the subscription. If
installation sets a client property, close restores or removes that property
according to the previous state. These rules sound mechanical because they are.
Mechanical lifecycle rules are the ones most worth encoding.

\api{helpOnF1} is a good restoration example. The behavior obtains the focused
input map, action map, F1 keystroke, previous input mapping, and previous action
for the generated action key. Installation writes the new mapping and action.
Cleanup restores the old mapping and old action. This is more careful than
calling \api{inputMap.remove} and \api{actionMap.remove} unconditionally,
because the component may have had a meaningful F1 action before the behavior
was installed. A reusable behavior should be a guest in the component, not an
owner of everything the component previously knew.

Tooltip composition is the example where restoration and cardinality meet. A
component has one tooltip text property. Validation, disabled reason, static
help, and F1 help indication all want to contribute. If each contributor is a
separate behavior with its own tooltip writer, restoration becomes impossible:
which behavior owns the original tooltip, and which behavior restores it last?
The composed tooltip behavior solves this by making the tooltip slot a single
owned behavior and treating validation and disabled reason as inputs. The
tooltip slot is single-cardinality; the information sources are multiple model
facts.

This is the pattern to copy for other single slots. A component has one border
property. It has one accessible name and one accessible description. It has
specific input-map keys. It has specific action-map keys. If several behaviors
need to influence one slot, either compose the inputs under one behavior or use
a deliberately designed composition object. Do not let unrelated behaviors
race to write the same Swing property. Last-writer-wins is not a lifecycle
policy.

Sometimes the correct composition boundary is visual rather than property-based.
Validation border, focus ring, dirty marker, and busy mark may all want to
decorate a component. Replacing borders for every mark can become brittle under
Look-and-Feel changes. A \api{JLayer}, parent overlay, or painting behavior may
be a better owner for visual marks that cross component boundaries or that
should not disturb the component's original border. The behavior system does
not force every decoration into a border; it only asks that whichever mechanism
is chosen has a key, phase, cardinality, and close contract.

Custom behaviors should be small enough to review in one sitting. If a behavior
needs to know the whole form, coordinates between several components, and a
presenter service, it may be a screen workflow rather than a reusable component
capability. A behavior that adds select-all-on-focus is reusable. A behavior
that opens a date chooser for one field may be reusable if it owns one
component and one action. A behavior that validates two fields, enables a
button, changes a table selection, and calls a repository is not a behavior; it
is presenter logic trying to hide as view assembly.

A useful test is the sentence "this component can..." A text field can select
all on focus. It can open help on F1. It can show a validation tooltip. It can
bind raw text to a model. It can publish status while focused. These are
component capabilities. "When the user changes credit limit, reload pricing
rules and update the table" is not a component capability. It is workflow. The
workflow may use commands, tasks, values, and bindings, but it should not be
installed as a tiny component behavior merely because it happens to start from
a field event.

Custom behavior factories should validate required collaborators early. If a
behavior needs a model value, command, resource text, or topic, the factory can
check nulls before returning the behavior object. If it needs a context service
such as help, installation can fail with a clear message when the scope lacks
that service. This split keeps construction failures near construction and
context failures near installation. A behavior that waits until the user presses
F1 to discover there is no help service has failed too late.

The behavior object itself should not usually be stateful across installations.
It may capture immutable collaborators such as a model, resource string,
command, topic, or policy. It should be careful with mutable per-component
state. The same behavior instance might be reused accidentally across
components, or a generated plan might create new instances for each component.
The safest design is for \api{install} to allocate per-component listener
state, snapshot previous component state, and return a binding that closes that
state. Captured mutable state should be intentional and documented.

Client properties are sometimes a pragmatic bridge. A behavior may set a
component client property for a renderer, overlay, test key, or diagnostic
marker. Client properties are ordinary Swing state, so the same restoration
rule applies. Snapshot whether the property existed and what value it had.
Close should restore the old value or remove the property if it was absent. A
behavior that leaves a diagnostic property behind can confuse later tests and
support tools.

Input maps and action maps need particular care because they are layered by
focus condition. A behavior that installs \api{WHEN\_FOCUSED} F1 help should
not accidentally modify \api{WHEN\_ANCESTOR\_OF\_FOCUSED\_COMPONENT}. A behavior
that installs Enter-to-commit in a text field should consider whether the text
component already uses Enter for another purpose. For single-line fields,
Enter often commits. For text areas, Enter inserts a newline unless the product
explicitly changes that behavior. The behavior descriptor cannot answer that
product question; the factory choice must.

Key behaviors should prefer input maps and action maps over low-level key
listeners when they represent commands. The built-in \api{commitOnEnter}
behavior uses a key listener because it is deliberately small and text-component
specific, but command-oriented behaviors should usually use Swing's key binding
system. Key bindings cooperate better with focus traversal, Look-and-Feel, and
action enablement. A custom behavior should choose the Swing mechanism that
matches the responsibility, not the one that is shortest to type.

Busy overlays show a different kind of custom behavior. The target is a
\api{JLayer}, not the wrapped view. The behavior observes a readable busy value
and asks the layer to repaint while active. It does not own the business task
that produces busy state. It does not decide whether commands are disabled. It
does not load data. It owns the visual overlay relationship between a busy fact
and a Swing layer. That narrowness is why the same behavior can be reused in
dialogs, tables, and full screens.

Command behaviors belong in the same family. A button, menu item, toolbar
item, or key binding adapts a command's identity, enabled state, selected state,
text, icon, tooltip, and handler to Swing. The command owns the operation. The
behavior owns the component relationship. This prevents the classic Swing
problem where one toolbar button is disabled, the menu item remains enabled,
and the key binding still fires. One command fact should drive every Swing
surface that presents it.

Custom behavior tests should be explicit about restoration. Record the initial
listener count, input-map value, action-map value, border, tooltip, accessible
text, or client property. Install the behavior. Assert the installed effect.
Close the scope. Assert the previous state. Then, when useful, mutate the model
after close and assert the component no longer follows. These tests are not
overly defensive. They document the lifecycle contract that makes the behavior
safe to generate and repeat.

The most common custom-behavior mistake is to return \api{Binding.noop} or an
empty close action after installing real state. That is acceptable only for a
behavior that truly installs no disposable or restorable state. If a behavior
adds a listener and returns no cleanup, it has turned a view activation into a
memory and behavior leak. If a behavior changes a component property and
returns no restoration, it has made close order irrelevant because close cannot
repair anything. The code may still pass a visual smoke test; it fails the
lifecycle contract.

Another mistake is to perform model mutation during installation. Installation
should adapt current model state to the component and install listeners for
future changes. It should not decide that a field is invalid, trigger a save,
open help, or run a task merely because a component exists. Those actions
belong to models, presenters, commands, or tasks. A behavior installation that
has side effects beyond component setup is difficult to test and dangerous in
generated plans because generation may install it earlier or more often than a
handwritten presenter once did.

The better mental model is that a behavior is a small adapter with a passport.
The adapter knows how to install one component capability. The passport is its
descriptor: key, phase, cardinality, and primary-binding conflict. The scope
checks passports, admits behaviors in order, owns their close handles, and can
report which passports are installed. This is enough structure to make Swing
assembly safe without pretending that Swing has stopped being Swing.

\section{Plain Swing baseline}

The plain Swing version of a field setup may look like this:

\begin{lstlisting}[caption={Handwritten field wiring.}]
nameField.getDocument().addDocumentListener(new DocumentListener() {
    @Override
    public void insertUpdate(DocumentEvent event) {
        form.name.setRawText(nameField.getText(), StandardChangeOrigin.USER);
    }
    @Override
    public void removeUpdate(DocumentEvent event) {
        form.name.setRawText(nameField.getText(), StandardChangeOrigin.USER);
    }
    @Override
    public void changedUpdate(DocumentEvent event) {
        form.name.setRawText(nameField.getText(), StandardChangeOrigin.USER);
    }
});

form.name.problemSet().subscribe(event ->
        nameField.setToolTipText(primaryProblemText(event.newValue())));
\end{lstlisting}

This code omits feedback guards, EDT assertions, close handles, tooltip
composition, border restoration, and conflict detection. A careful developer can
add all of them. Corusco's value is that the careful version becomes the
ordinary version.

\textbf{From plain Swing to Corusco.}

Use direct bindings when there are only one or two relationships and a generated
plan would be noise. Use behaviors when capabilities repeat, when generated
companions participate, or when conflict detection matters. Use a scope in both
cases.

The decision is not "framework or raw Swing". It is "which relationships are
repeated enough to deserve metadata and a plan?"

Adoption usually works best in layers. Start by replacing one ad hoc document
listener with a direct binding and putting that binding into a scope. Then move
validation tooltip and border logic into separate bindings. When the same
sequence appears on several fields, introduce behaviors or generated behavior
plans. This keeps the migration understandable: each step gives one existing
relationship a clearer owner. A large screen does not need to jump from raw
Swing to generated plans in one commit.

The important discipline is that the old listener must actually disappear. If a
component is bound through Corusco and still has an old listener that writes to
the same model, ownership is no clearer than before. During review, look for
"parallel wiring": a binding installed in one method and a listener installed in
another method, both trying to keep the same fact synchronized. That is a sign
that the migration stopped halfway.

\textbf{Testing bindings and behaviors.}

Binding tests should run on the EDT. They should assert initial component state,
component-to-model propagation, model-to-component propagation, and cleanup.
Behavior tests should also assert installed behavior keys, ordering when phases
matter, and conflict detection.

\begin{lstlisting}[caption={Behavior conflict test shape.}]
BehaviorScope scope = new BehaviorScope();

scope.install(nameField, List.of(
        StandardBehaviors.textFieldBinding(form.name)));

assertThatThrownBy(() -> scope.install(nameField, List.of(
        StandardBehaviors.textFieldBinding(form.nickname))))
        .isInstanceOf(BehaviorConflictException.class);
\end{lstlisting}

\textbf{Worked example.}

\begin{examplebox}{Bindings}
\swingIntegrationExample{} contains direct bindings and behavior-style
installation in compact form. Read it by following lifetimes: model lifetime,
component lifetime, binding lifetime, and scope lifetime.
\end{examplebox}

\textbf{Review notes.}

Bindings are attractive because they remove listener code from the view. That
does not mean every relationship should be hidden. Use names and descriptors
that make the installed plan inspectable. A generated plan should be boring to
read, not magical.

Good binding code should be inspectable at two levels. At the local level, a
reader should be able to see that \api{nameField} is bound to \api{form.name},
that validation problems decorate the same component, and that the scope closes
with the panel. At the generated-plan level, a reader should be able to open the
generated behavior plan and recognize the same facts: primary text binding,
tooltip, border, help, accessibility, and interaction behavior. If generated
code hides the relationship so well that the reader cannot reconstruct it, the
plan needs better names or smaller generated units.

The simplest binding is a promise between two mutable places: when one changes,
the other is brought into agreement. That promise is easy to say and surprisingly
easy to get wrong. Initial synchronization must choose a direction. User edits
must not be mistaken for model resets. Programmatic model changes must not loop
back as fresh user edits. Cleanup must remove every listener installed to keep
the promise. A handwritten listener can satisfy all of these rules, but the
rules are usually scattered.

Corusco bindings put the promise into one object. The object can assert the EDT
boundary, install component and model listeners together, apply feedback guards,
and return a close handle. This is not meant to make Swing invisible. It is meant
to make synchronization reviewable. A binding whose source, target, direction,
origin policy, and lifetime are visible is easier to trust than three anonymous
listeners that happen to cooperate.

Behaviors become useful when a component has several independent promises. A
text field may have a primary text binding, validation tooltip, validation
border, dirty marker, help key, accessibility name, and focus-selection rule.
Those promises should not all live in one large "bind everything" method. They
also should not be installed in random order by unrelated code. A behavior plan
keeps them separate while still letting the screen install them as a coherent
unit.

Phases prevent accidental ordering from becoming policy. A primary value binding
should be established before decoration that depends on the field model. A
tooltip composer should run after contributors have been registered. A focus
behavior should not override a validation behavior merely because its method was
called later. Naming phases turns order from folklore into data that generated
plans and tests can inspect.

Cardinality is the other half of conflict detection. A text field can have one
primary text binding, but several tooltip contributors. A component can have one
accessible-name owner, but several optional visual decorations. If all behaviors
are treated as independent, duplicate primary bindings silently fight. If all
behaviors are treated as exclusive, useful composition becomes impossible. The
behavior descriptor should say which kind of relationship it installs.

Generated behavior plans are most valuable when they are dull. A generated form
view contract says which components exist. The generated descriptors say which
field each component represents. The behavior plan says which standard
relationships should be installed. The handwritten view still chooses geometry
and component construction. The generated plan merely prevents every screen from
rewriting the same binding recipe by hand.

Custom behaviors should be small and reversible. A behavior that installs a key
binding should remove that key binding. A behavior that changes a border should
restore the previous border or cooperate with a composed border policy. A
behavior that listens to a model should close its subscription. If a custom
behavior needs five services and several unrelated components, it is probably a
presenter workflow, not a reusable view behavior.

There is a useful migration smell: a method named \api{wireListeners}. It often
contains model synchronization, command enablement, tooltip updates, validation
decoration, focus behavior, and cleanup omissions in one place. Do not replace
the whole method at once. Extract one promise, give it a binding or behavior,
and close it through a scope. Repeat until the method either disappears or
contains only genuinely local listener work.

The behavior scope is not a dependency injection container. It should know what
relationships are installed on which components and how to close them. It should
not become the place where repositories, presenters, resources, and every
application service are discovered. Keep the context narrow. Pass only the
metadata and services a behavior needs to install its component relationship.

The most useful tests are deliberately mechanical. Install the plan. Verify the
component has the expected value. Change the component and observe the model.
Change the model and observe the component. Close the scope and verify later
changes do not cross the boundary. Add a duplicate primary behavior and assert a
conflict. These tests may look repetitive, but they defend the invariants that
make generated plans safe to use.

\section{Migration, testing, and review scenarios}

A good migration from listener-heavy Swing to Corusco bindings rarely begins
with generation. It begins with naming one relationship. Find a document
listener that writes field text to a model. Replace it with a direct binding.
Put the binding in a scope. Close the scope when the panel or dialog closes.
Run the screen. Then find the tooltip listener that recomputes validation and
replace it with a problem-set tooltip binding. Then replace the border
decoration. This order is intentionally conservative: each step removes one
unowned relationship and gives it an explicit lifecycle.

The migration should be reviewed by looking for remaining duplicate ownership.
If the same field is bound to a model and also has an old listener that writes
the model, the screen has not improved. If a generated plan installs tooltip
behavior and a presenter later calls \api{setToolTipText} directly, the tooltip
slot has two owners. If a behavior installs F1 help and old key-listener code
also opens help, key routing has two owners. Corusco makes ownership visible,
but it cannot make code outside the scope disappear. The reviewer must look for
parallel wiring and remove it deliberately.

One useful migration artifact is a relationship inventory. For each important
component, list the model fact it presents, validation facts it decorates,
commands it invokes, focus behavior it owns, help topic it opens, accessible
metadata it exposes, and cleanup owner. The inventory need not be a permanent
document. It is a way to discover that one field has eight relationships hidden
in three methods and two anonymous listeners. After migration, the same
inventory should map naturally to direct bindings or behaviors in one scope.

Generated plans should be introduced when the inventory repeats. A single
custom field does not justify generation. Ten text fields that all use the same
form-model pattern do. The generated plan should be compared with the manual
inventory: text binding, validation tooltip, validation border, select all on
focus, accessible text when metadata is available, help when topic metadata is
available. If generation omits a capability, either the generator should grow
or the handwritten extension should add the capability through the same scope.
Do not patch generated omissions with unscopeable listeners.

Testing follows the same layers. Core model tests prove that raw text, parsed
value, problems, dirty state, and command enablement behave without Swing.
Direct binding tests prove that a Swing component mirrors one model fact and
cleans up. Behavior tests prove that a reusable capability installs and
restores its Swing state. Generated-plan tests prove that the right behaviors
are installed for generated view contracts. Screenshot tests prove that the
assembled screen looks coherent. These layers are complementary; none should
be asked to prove everything.

A direct binding test should be almost boring. Construct a model and component
on the EDT. Install the binding. Assert initial component state. Perform a
user-like component change. Assert the model and, when relevant, the origin or
problem state. Perform a model change. Assert the component. Close. Change the
component again. Assert the model no longer follows. The value of the test is
not the number of lines saved. The value is that the relationship's lifetime is
executable documentation.

A behavior test adds descriptor checks. Install a behavior list in an
intentionally scrambled order and assert that \api{installedBehaviorKeys}
reports binding, decoration, and interaction phases in the expected order.
Install two tooltip behaviors and assert a duplicate-key conflict. Install two
primary bindings with different keys and assert a primary-binding conflict.
Install a help behavior without a help service and assert the failure message.
These tests are close to the source because they are testing assembly rules,
not business meaning.

Generated-plan tests should avoid overfitting to generated formatting. The
source generator can change indentation or helper names without changing the
plan contract. Tests can compile a sample annotated record, call the generated
plan, and inspect behavior keys or effects on components. Processor tests may
also assert that generated source contains important calls when those calls are
part of the public generation contract. Use source-text assertions sparingly
and prefer compiled behavior when possible.

The most useful screen-level test is often a hybrid. Use \api{SwingMvpTester}
or an EDT harness to construct the view and presenter. Query the behavior scope
for expected behavior keys on named components. Then interact with the
components and assert model state. Finally close the presenter or view and
assert that behavior keys are gone or that model changes no longer reach the
components. This test proves the plan is installed, observable, and disposable
in the same lifecycle that the real screen uses.

Screenshot tests belong after these tests, not instead of them. A screenshot
can reveal a missing validation border, a stale tooltip policy visible in a
hover capture, a busy overlay that covers the wrong component, or a command
button that remains enabled. It cannot reliably prove that cleanup removed a
document listener. It cannot prove that user origin is preserved. It should be
the visual smoke layer on top of model and binding tests.

Code review can use a small checklist. Does every installed relationship have
a close owner? Does the scope live for exactly one view or dialog activation?
Are bindings installed and closed on the EDT? Is there only one primary binding
per component? Does tooltip ownership go through a composed tooltip when
several facts contribute? Does a custom behavior restore what it replaces? Do
generated and handwritten behaviors enter the same scope? Are old direct
listeners removed after migration? The checklist is short because the model is
simple: relationships must be named, ordered, conflict-checked, and closed.

There are also signs that a binding abstraction has gone too far. A behavior
that hides significant business rules is a problem. A generated plan that
builds layout is a problem when the application needs MigLayout-grade control.
A binding that accepts a whole presenter and calls arbitrary methods is a
problem because it becomes impossible to know what relationship it owns. A
scope that is stored globally is a problem because scopes are lifecycles, not
registries. Corusco's Swing integration is intentionally modest: it structures
component relationships; it does not replace presenter design.

The converse problem is refusing abstraction after the pattern is obvious. A
form with twenty fields should not hand-write the same text binding, tooltip,
border, and focus behavior twenty times. A screen with several buttons and menu
items should not manually copy command enabled state into each surface. A table
screen should not manually synchronize selection indexes to model ids in three
places. When the same relationship pattern repeats, a behavior, descriptor, or
generated plan is not ornament. It is the mechanism that keeps repeated Swing
machinery consistent.

The final migration habit is to keep examples executable. A chapter listing
that says "install generated behavior plan here" is useful only if an example
actually compiles through the annotation processor or mirrors the generated
plan for a documented reason. The companion examples should contain at least
one direct binding example, one behavior-scope example, one generated plan
example, and tests that close the scopes. This keeps the book honest: every
claimed simplification must survive javac and an EDT test.

\section{Exercises}

\begin{exercisebox}
\begin{enumerate}
\item Replace one document listener with \api{BindingFactory.textField} and add
a cleanup test.
\item Split a field setup into primary binding, validation tooltip, validation
border, and focus status behavior.
\item Install two primary bindings on one component in a test and verify the
conflict.
\item Write a custom reversible behavior that installs a key binding and
removes it on close.
\item Decide whether a small one-off panel should use direct bindings or a
behavior plan. Write down the reason.
\item Add a behavior-scope test that inspects installed behavior keys for one
generated field.
\end{enumerate}
\end{exercisebox}

\textbf{Checklist.}

\begin{itemize}
\item Use direct bindings for one simple model/component relationship.
\item Use behaviors when capabilities repeat or generated plans must cooperate.
\item Install and close bindings on the EDT.
\item Close scopes with the owning view or dialog.
\item Keep primary value binding separate from validation and tooltip
decoration.
\item Verify conflict detection for primary bindings and single-cardinality
behaviors.
\item Write custom behaviors as small reversible units.
\item Avoid turning behavior scopes into dependency injection containers.
\end{itemize}
